%%vsgtu
%
% %%%

%%%   Самарский государственный технический университет

%%%   Кафедра прикладной математики и информатики

%%%

%%%   Преамбула для статьи в журнал "Вестник Самарского государственного

%%%   технического университета. Серия: «Физико-математические науки»"

%%%   Ориентирована под MikTEx 2.4 for Win32

%%%

%%%   Хотя сам журнал собирается под GNU/Linux на дистрибутиве TeXLive



\documentclass[11pt,a4paper,twoside]{article}



\usepackage[cp1251]{inputenc}

\usepackage[T2A]{fontenc}

\usepackage[english,russian]{babel}

\usepackage{samgtubib}

\usepackage[dvips]{graphicx}

\usepackage[multiple]{footmisc}



\usepackage{samgtu}

\renewcommand{\VestnikYear}{2009} % Год издания Вестника

\renewcommand{\VestnikNumber}{2\,(19)} % Номер Вестника

\renewcommand{\VestnikMonth}{Ноябрь} % Месяц выхода Вестника

\renewcommand{\VestnikData}{01.11.09} % Дата подписания Вестника в печать

\renewcommand{\VestnikUPL}{26,64} % Усл. печ. л.

\renewcommand{\VestnikUIL}{25,73} % Уч.-изд. л.

\renewcommand{\VestnikRegNumber}{479/09} % Регистрационный номер

\renewcommand{\VestnikOrder}{994} % Номер заказа






%
% \allowdisplaybreaks[4]
% 
% \begin{document} 

 
% \tableofcontents
% \newpage

\renewcommand{\firstpage}{190}
\renewcommand{\lastpage}{203}


%%%%%%%%%%%%%%%%%%%%%%%%%%%%%%%%%%%%%%%%%%%%%%%%%%%%%%%%%%%%%%%%%%%%%%%%%%%%%%%%%%%%
%Информация о статье и об авторах на русском и английском языках

% Номер УДК
\udk{539.3}
\msc{74B05, 74S20}

\title{Несвязанная осесимметричная динамическая задача  обратного пьезоэффекта для радиально  поляризованного цилиндра }% Полный заголовок
{Несвязанная осесимметричная динамическая задача \dots}% Заголовок, который идёт в верхний колонтитул

\author{Д.\,А.~Шляхин}
{Ш\,л\,я\,х\,и\,н~Д.\,А.}% Автор(ы) для оглавления и колонтитула через \,

\institute{\sgasu.}


\email{E-mail}{d-612-mit2009@yandex.ru} 

\otherinf{\fullauthor{Дмитрий Аверкиевич Шляхин} \regalia{к.т.н., доц.} \position{докторант}
\dept{каф. сопротивления материалов и строительной механики} }


\keywords{задача обратного пьезоэффекта, пьезокерамический цилиндр, осесимметричная динамическая нагрузка}

\workabstract{Рассматривается несвязанная осесимметричная нестационарная задача обратного пьезоэффекта для анизотропного пьезокерамического радиально поляризованного  цилиндра при действии на его внешних радиальных поверхностях электрического потенциала, являющегося произвольной функцией по аксиальной координате и~времени. 
Новое замкнутое решение построено методом разложения по собственным вектор"=функциям в~форме структурного алгоритма конечных преобразований. 
Полученные соотношения позволяют определять частоты собственных колебаний, напряжённо"=деформированное состояние элемента, а~также потенциал и~напряжённость индуцируемого электрического поля.}

\engtitle{Untied Axisymmetris Dynamic Problem of the Return Piezoeffect for the Radially Polarised Cylinder}

\engauthor{D.\,A.~Shljakhin}{S\,h\,l\,j\,a\,k\,h\,i\,n~D.\,A.}


\enginstitute{\engsgasu.}

\otherenginfm{\fullauthor{Dmitry  A. Shljakhin} \regalia{Ph. D. (Techn.)} \position{Doctoral Candidate}
\dept{Dept. of Resistance of Materials \& Construction Mechanics}}

\engabstract{The non-stationary problem of the return piesoeffect for anisotropic piesoceramic radially polarized cylinder is considered untied axisymmetris at action on its external radial surfaces of the electric potential which is any function on axial co-ordinate and time. The new closed decision is developed by a decomposition method to its own  vector--functions in the form of structural algorithm of final transformations. The received ratios allow to define frequencies of own fluctuations, the is intense-deformed condition of an element, and also potential and intensity of induced electric field.}

\engkeywords{problem of the return piezoeffect, the piezoceramic cylinder, axisymmetris dynamic loading}

\date{12/X/2009} % Дата поступления статьи в редакцию.
\revisiondate{29/IX/2010} % В окончательном виде

%%%%%%%%%%%%%%%%%%%%%%%%%%%%%%%%%%%%%%%%%%%%%%%%%%%%%%%%%%%%%%%%%%%%%%%%%%%%%%%%%%%%



\maketitle






\Section[N]{Введение} Исследование вынужденных осесимметричных колебаний пьезокерамического цилиндра при радиальной поляризации материала приводит к~формированию сложной системы дифференциальных уравнений. 
Проблемы интегрирования расчётных соотношений решаются путём анализа более простых задач электроупругости о~колебаниях бесконечного цилиндра~[\citen{sl:c:1, sl:c:2}], исследования тонких пластин и~оболочек [\citen{sl:c:1, sl:c:2, sl:c:3}].
Кроме того, используются численные методы расчёта~[\citen{sl:c:2, sl:c:3}].

В настоящей работе приводится новое замкнутое решение осесимметричной  несвязанной динамической задачи обратного пьезоэффекта в~трёхмерной постановке. Построенный для пьезокерамического цилиндра алгоритм позволяет удовлетворить смешанные краевые условия. 
Данный подход был предложен и~впервые использован в~работах Ю.~Э.~Сеницкого при исследовании короткого анизотропного цилиндра~[\citen{sl:c:4, sl:c:5, sl:c:6}].

\Section{Постановка задачи} 
В~работе исследуется полый анизотропный цилиндр с~наведённой радиальной поляризацией, выполненный из пьезокерамического материала, и~занимающий в~цилиндрической системе  координат $(r_* , \theta, z_*  )$ область $\Omega $: 
$\{ a\le r_*  \le b, 0\le \theta \le 2\pi, 0\le z_*  \le h\}$. 

Для рассматриваемой задачи можно сформулировать различные физически реализуемые краевые условия. 
Для определённости принимаем неэлектродированные торцевые плоскости свободными от механических напряжений, а~радиальные поверхности "--- электродированными с~заземлением закреплённой ее внутренней части.

Краевая задача моделирует работу пьезоэлементов в~приборах обратного пьезоэффекта при действии на внешней криволинейной поверхности цилиндра электрического потенциала $V^*  ( z_* , t_*  )$. 
В~результате данного воздействия исследуемый элемент деформируется.

В~общем случае дифференциальные уравнения осесимметричного движения и~электростатики однородной упругой анизотропной среды в~цилиндрической 
системе координат записываются в~следующем виде~[\citen{sl:c:2, sl:c:3}]:
\begin{equation}
\label{sl:eq1}
\begin{array}{r}
\displaystyle
\frac{\partial \sigma_{rr} }{\partial r_*  }+\frac{\partial \sigma _{rz} }{\partial z_*  }+\frac{\sigma _{rr} -\sigma _{\theta \theta } }{r_*  
}-\rho \frac{\partial ^2U^*  }{\partial t_*  ^2 }=0, \vspace{1mm}
\\
\displaystyle
\frac{\partial \sigma _{rz} }{\partial r_*  }+\frac{\partial \sigma _{zz} }{\partial z_*  }+\frac{\sigma _{rz} }{r_*  }-\rho \frac{\partial 
^2W^*  }{\partial t_*  ^2 }=0, \vspace{1mm}
\\
\displaystyle
\frac{\partial D_r }{\partial r_*  }+\frac{D_r }{r_*  }+\frac{\partial 
D_z }{\partial z_*  }=0.
\end{array}
\end{equation}

При радиальной поляризации материала уравнения состояния пьезокерамического тела определяются следующими равенствами~[\citen{sl:c:2, sl:c:3}]:
\begin{equation}
\label{sl:eq2}
\begin{array}{l}
\displaystyle
\sigma _{rr} =C_{33} \frac{\partial U^*  }{\partial r_*  }+C_{13} \Bigl( 
\frac{U^*  }{r_*  }+\frac{\partial W^*  }{\partial z_*  }
\Bigr)+e_{33} \frac{\partial \phi ^*  }{\partial r_*  }, \vspace{1mm}
\\
\displaystyle
\sigma _{\theta \theta } =C_{13} \frac{\partial U^*  }{\partial r_*  
}+C_{11} \frac{U^*  }{r_*  }+C_{12} \frac{\partial W^*  }{\partial 
z_*  }+e_{31} \frac{\partial \phi ^*  }{\partial r_*  }, \vspace{1mm}
\\
\displaystyle
\sigma _{zz} =C_{13} \frac{\partial U^*  }{\partial r_*  }+C_{12} 
\frac{U^*  }{r_*  }+C_{11} \frac{\partial W^*  }{\partial z_*  
}+e_{31} \frac{\partial \phi ^*  }{\partial r_*  }, \vspace{1mm}
\\
\displaystyle
\sigma _{rz} =C_{55} \Bigl( \frac{\partial W^*  }{\partial r_*  
}+\frac{\partial U^*  }{\partial z_*  } \Bigr)+e_{15} \frac{\partial 
\phi ^*  }{\partial z_*  };
\end{array}
\end{equation}
\begin{equation}
\label{sl:eq3}
\begin{array}{l}
\displaystyle
D_r =-\varepsilon _{33} \frac{\partial \phi ^*  }{\partial r_*  }+e_{31} 
\Bigl( \frac{U^*  }{r_*  }+\frac{\partial W^*  }{\partial z_*  } \Bigr)+e_{33} \frac{\partial U^*  }{\partial r_*  }, \vspace{1mm}
\\
\displaystyle
D_z =-\varepsilon _{11} \frac{\partial \phi ^*  }{\partial z_*  }+e_{15} 
\Bigl(\frac{\partial W^*  }{\partial r_*  }+\frac{\partial U^*   }{\partial z_*  } \Bigr).
\end{array}
\end{equation}
Здесь $\sigma _{jk} (r_* , z_* , t_*  )$, $U^* (r_* , z_* , t_* )$, $W^*  (r_* , z_* , t_* )$ "--- соответственно компоненты тензора 
механических напряжений и вектора перемещений, $j$, $k\in \{r, \theta, z \}$; 
$D_r (r_*, z_*, t_* )$, $D_z(r_*, z_*, t_*)$, $\phi ^*  (r_*  , z_*  , t_*  )$ "--- компоненты вектора индукции и потенциал электрического поля;
$\rho$, $C_{ms}$, $e_{ms}$ "--- объёмная плотность, модули упругости и пьезомодули анизотропного электроупругого материала,
$m$, $s\in\{1, 2,$ $3, 4, 5\}$; $\varepsilon _{11}$, $\varepsilon _{33} $ "--- диэлектрические проницаемости.

При решении несвязанной задачи обратного пьезоэффекта используется допущение, что механические деформации не оказывают влияние на электрическое поле. 
Таким образом, в равенствах (\ref{sl:eq3}) следует исключить компоненты вектора перемещений: 
\begin{equation}
\label{sl:eq4}
U^*  =W^*  =0.
\end{equation}
В результате подстановки (\ref{sl:eq3}) в  (\ref{sl:eq1}) с учётом (\ref{sl:eq4}) получаем дифференциальное уравнение и~краевые условия для электрического потенциала:
\begin{equation}
\label{sl:eq5}
\nabla ^2\phi +\frac{\varepsilon _{11} }{\varepsilon _{33} }\frac{\partial ^2\phi }{\partial z^2}=0,
\end{equation}
где $\nabla^2 = \partial^2 / \partial^2 r +r^{-1} \partial / \partial r$;
при $z=0$ и $z=L$ имеем
\begin{equation}
\label{sl:eq6}
D_z =-C_{33} \varepsilon _{11} e_{15}^{-1} \frac{\partial \phi }{\partial z}=0,
\end{equation}
а при $r=k$ и $r=1$ "---
\begin{equation}
\label{sl:eq7}
\phi ( k, z, t)=0, \quad \phi (1, z, t)=V(z, t). 
\end{equation}
Здесь используются следующие безразмерные величины:
 $r=r_*/b$, $z=z_*/b$, $L=h/b$, $k=a/b < 1$; $t= t_* \sqrt{C_{33}/\rho}/b$;
$\phi=\phi^* e_{15}/(bC_{33})$, $V=V^* e_{15}/(bC_{33})$.


На втором этапе исследования рассматривается начально"=краевая задача теории  упругости в предположении, что потенциал электрического поля $\phi$
известен (найден).

В результате подстановки равенств (\ref{sl:eq2}) в (\ref{sl:eq1}) получаем систему дифференциальных уравнений, граничные и~начальные условия для динамической задачи теории упругости в~безразмерной форме:
\begin{equation}
\label{sl:eq8} \hspace{-1mm}
\begin{array}{c}
\displaystyle
\nabla ^2U-\frac{C_{11} }{C_{33} }\frac{U}{r^2}+\frac{C_{55} }{C_{33}}\frac{\partial ^2U}{\partial z^2}+\frac{C_{13} + C_{55}}{C_{33} }
\frac{\partial ^2W}{\partial r\partial z}+\frac{C_{13} -C_{12}}{C_{33} }\frac{1}{r}\frac{\partial W}{\partial z}-
\frac{\partial ^2U}{\partial t^2}=B_1 , \vspace{2mm}
\\ 
\displaystyle
\frac{C_{55} }{C_{33} }\nabla ^2W+\frac{C_{11} }{C_{33} }\frac{\partial ^2W}{\partial z^2}+\frac{C_{13} +C_{55} }{C_{33} }
\frac{\partial ^2U}{\partial r\partial z}+\frac{C_{12} +C_{55}}{C_{33} }\frac{1}{r}\frac{\partial U}{\partial z}-\frac{\partial ^2W}{\partial t^2}=B_2 , 
\end{array} \hspace{-6mm}
\end{equation}
где
$$B_1 =-\frac{e_{33} }{e_{15} }\nabla ^2\phi +\frac{e_{31} }{e_{15} 
}\frac{1}{r}\frac{\partial \phi }{\partial r}-\frac{\partial ^2\phi 
}{\partial z^2} ,\quad  B_2 =-\frac{e_{15} +e_{31} }{e_{15} }\frac{\partial 
^2\phi }{\partial r\partial z}-\frac{1}{r}\frac{\partial \phi }{\partial 
z};$$
при $z=0$ и $z=L$ имеем
\begin{equation}
\label{sl:eq9}
\begin{array}{c}
\displaystyle
\sigma _{zz} =C_{13} \frac{\partial U}{\partial r}+C_{12} \frac{U}{r}+C_{11} \frac{\partial W}{\partial z}+C_{33} \frac{e_{31} }{e_{15} }\frac{\partial \phi }{\partial r}=0, \vspace{1mm}
\\
U(r, 0, t)=U_1 (r,  t), \quad U(r, L,  t)=U_2 (r,  t);
\end{array}
\end{equation}
при $r=k$ и $r= 1$ "---
\begin{equation}
\label{sl:eq10}
\begin{array}{c}
U(k, z,  t)=W(k, z,  t)=0, \vspace{1mm}
\\
\displaystyle
\sigma _{rr} \big|_{r=1} =C_{33} \frac{\partial U}{\partial r}+C_{13} \Bigl(U+\frac{\partial W}{\partial z} \Bigr)+
C_{33} \frac{e_{33} }{e_{15}}\frac{\partial \phi }{\partial r}=0, 
\vspace{1mm}
\\
\displaystyle
\sigma _{rz}\big|_{r=1} =C_{55} \Bigl(\frac{\partial W}{\partial r}+\frac{\partial U}{\partial z} \Bigr)+C_{33} \frac{\partial \phi }{\partial z}=0;
\end{array}
\end{equation}
при $t=0$ "---
\begin{equation}
\label{sl:eq11}
\begin{array}{c}
U(r, z, 0)=U_0 (r, z),\quad \dot {U}(r, z, 0)=\dot {U}_0 (r, z);
\\
W(r, z, 0)=W_0 (r, z),\quad \dot {W}(r, z, 0)=\dot {W}_0 (r, z).
\end{array}
\end{equation}
Здесь используются следующие безразмерные величины:
$U=U^*/b$, $V=V^*/b$, $U_1=U_1^*/b$, $U_2=U_2^*/b$, $U_0=U_0^*/b$; точка означает дифференцирование по $t$;
$U_1^*$, $U_2^*$  "--- известные радиальные перемещения на торцевых поверхностях; 
$U_0^*$, $\dot {U}_0^*$, $W_0^*$, $\dot {W}_0^*$ "--- известные в~начальный момент времени перемещения, скорости перемещений.

Соотношения \eqref{sl:eq5}--(\ref{sl:eq11}) представляют собой  математическую формулировку рассматриваемой несвязанной начально"=краевой задачи электроупругости. 

\smallskip
\Section{Построение общего решения краевой задачи электростатики}
Решение осуществляется методом интегральных преобразований с последовательным использованием конечного косинус"=преобразования Фурье по координате $z$ и~интегрального преобразования Вебера [\citen{sl:c:7}] по переменной $r$.

Применяя к (\ref{sl:eq5})--(\ref{sl:eq7}) косинус"=преобразование Фурье по переменной $z$, получаем следующую задачу в пространстве изображений: 
\begin{equation}
\label{sl:eq12}
\nabla ^2\phi _c -\frac{\varepsilon _{11} }{\varepsilon _{33} }j_m^2 \phi _c =0, 
\end{equation}
где
$\phi _c ( r, m, t)= {\cal F}^m_c\bigl( \phi ( r, z, t) \bigr)$, $V_c (m, t) = {\cal F}^m_c \bigl(V_c(z, t)\bigr)$;
$$  
{\cal F}^m_c\bigl( {}\cdot {} \bigr) =
\int_0^L \bigl( {}\cdot {} \bigr)\cos j_m z\,  dz ,\quad j_m ={m\pi }/{L}; \quad m\in \{0\}\cup \mathbb N;
$$
при $r=k$ и $r= 1$ имеем
\begin{equation}
\label{sl:eq13}
\phi _c (k, m, t)=0, \quad \phi _c (1, m, t)=V_c (m, t).
\end{equation}

Используя процедуру стандартизации (приведение к однородным граничные условия) по координате $r$ к~краевой задаче (\ref{sl:eq12}), (\ref{sl:eq13}), представляем трансформанту Фурье в~виде 
\begin{equation}
\label{sl:eq14}
\phi _c (r, m, t)=(1-k)^{-1}(r-k)V_c +\chi _c (r, m, t)
\end{equation}
и получаем новую краевую задачу относительно $\chi _c (r, m, t)$:
\begin{equation}
\label{sl:eq15}
\nabla ^2\chi _c -\frac{\varepsilon _{11} }{\varepsilon _{33} }j_m^2 \chi _c =E_c , 
\end{equation}
где $E_c =\bigl[ \varepsilon _{11} \varepsilon _{33}^{-1} j_m^2 (r-k)-r^{-1}\bigr] V_c (1-k)^{-1}$; при $r=k$ и $r= 1$ имеем
\begin{equation}
\label{sl:eq16}
\chi _c (k, m, t)=\chi _c (1, m, t)=0.
\end{equation}

Применяем теперь к задаче (\ref{sl:eq15}), (\ref{sl:eq16}) конечное интегральное преобразование Вебера [\citen{sl:c:7}] по переменной $r$, вводя  трансформанту
\begin{equation}
\label{sl:eq17}
\psi _H (\xi _\kappa , m, t)=
\int_k^1 \chi_c (r, m, t) \bigl[Y_0 (\xi _\kappa k) J_0(\xi _\kappa r) - J_0 (\xi _\kappa k)Y_0 (\xi _\kappa r)\bigr]r\,dr 
\end{equation}
с соответствующей формулой обращения
\begin{equation}
\label{sl:eq18}
\chi _c (r, m, t)=2\pi^2\sum\limits_{\kappa=1}^\infty 
\frac{\psi_H ( \xi _\kappa , m, t)\bigl[Y_0( \xi _\kappa k) J_0 (\xi _\kappa r)-J_0 ( \xi _\kappa k)Y_0(\xi _\kappa r) \bigr]}
{\pi ^2\bigl[Y_0 ( \xi _\kappa k)J_0 (\xi _\kappa )-J_0 ( \xi _\kappa k)Y_0 ( \xi _\kappa )\bigr]^2-4\xi _\kappa^{-2} }.
\end{equation}
Здесь $\xi _\kappa$, $\kappa\in \mathbb N$ "--- положительные  параметры, определяемые из трансцендентного уравнения
\begin{equation}
\label{sl:eq19}
J_0 ( \xi _\kappa )Y_0 (\xi _\kappa k)-J_0 ( \xi _\kappa ) Y_0 (\xi _\kappa )=0.
\end{equation}
В результате получаем следующее выражение для трансформанты:
\begin{multline}
\label{sl:eq20}
\psi _H (\xi _\kappa , m, t)=-\Bigl(\xi _\kappa^2 +\frac{\varepsilon _{11} }{\varepsilon _{33} }j_m^2  \Bigr)^{-1} \times \\ \times 
\int_k^1 E_c \bigl[ Y_0 (\xi _\kappa k)J_0 (\xi _\kappa r)-J_0 (\xi _x k)Y_0 (\xi _\kappa r) \bigr]  r\, dr.
\end{multline}

Используя последовательно формулы обращения (\ref{sl:eq18}) и косинус"=преобразования Фурье, с учётом (\ref{sl:eq14}), (\ref{sl:eq20}) получаем выражение для $\phi (r, z, t)$:
\begin{multline}
\label{sl:eq21}
\phi (r, z, t)=\sum\limits_{m=0}^\infty 
\Omega _m^{-1} \cos (j_m z)\biggl\{ (1-k)^{-1}(r-k)V_c + \\+
2\pi^2\sum\limits_{\kappa=1}^\infty 
\frac{\psi_H ( \xi _\kappa , m, t)\bigl[Y_0( \xi _\kappa k) J_0 (\xi _\kappa r)-J_0 ( \xi _\kappa k)Y_0(\xi _\kappa r) \bigr]}
{\pi ^2\bigl[Y_0 ( \xi _\kappa k)J_0 (\xi _\kappa )-J_0 ( \xi _\kappa k)Y_0 ( \xi _\kappa )\bigr]^2-4\xi _\kappa^{-2} } \biggr\},
\end{multline}
где $ \Omega _m =\begin{cases}
 ~L, \text{ ~\,\,при $m=0$,}\\
 L/2, \text{ при $m\ne0$.} \end{cases}$





\Section{Построение общего решения начально"=краевой задачи теории упругости}
Решение осуществляется методом интегральных преобразований при последовательном использовании синус- и косинус"=преобразований Фурье с~конечными 
пределами по переменной $z$ [\citen{sl:c:7}] и~обобщённого конечного интегрального преобразование (КИП) [\citen{sl:c:6}] по радиальной координате $r$.

Приведём краевую задачу (\ref{sl:eq8})--(\ref{sl:eq11}) к~стандартной форме на основе следующего представления:
\begin{equation}
\label{sl:eq22}
U(r, z, t)=H_1 (r, z, t)+u(r, z, t), \quad W(r, z, t)=H_2 (r, z, t)+w(r, z, t), 
\end{equation}
где $$H_1 (r, z, t)=f_1 (z)U_1 (r, t)+f_2 (z)U_2 (r,  t),$$ 
\begin{multline*}
H_2 (r, z, t)=f_3 (z)\Bigl(U_1 ' +\frac{C_{12} }{C_{13} }r^{-1}U_1 \Bigr)+ \\+
f_4 (z)\Bigl(U_2 ' +\frac{C_{12} }{C_{13} }r^{-1}U_2  \Bigr)+
f_5 (z)\frac{\partial \phi }{\partial r}\Big|_{z=0} +f_6 (z)\frac{\partial \phi }{\partial r}\Big|_{z=L} .
\end{multline*}
Здесь и ниже <<штрих>> означает дифференцирование по соответствующей координате.

В результате подстановки (\ref{sl:eq22}) в (\ref{sl:eq8})--(\ref{sl:eq11}) с учётом условий 
\begin{equation}
\label{sl:eq23}
\begin{array}{c}
f_1 (L)=f_2 (0)=f_3 ' (L)=f_4 ' (0)=0, \qquad f_1 (0)=f_2 (L)=1,  \\ 
f_3 (0)=-f_3 (L), \qquad f_3 ' (0)=f_4 ' (L)=-C_{13} C_{11}^{-1}, \\ 
f_4 (0)=-f_4 (L), \qquad f_5 ' (0)=f_6 ' (L)=-C_{33} e_{31} (C_{11} e_{15} )^{-1}, \\
f_5 ' (L)=f_6 ' (0)=0, \qquad  f_5 (0)=-f_5 (L),  \qquad f_6 (0)=-f_6 (L)
\end{array}
\end{equation}
получаем новую краевую задачу относительно функций $u(r, z, t)$, $w(r, z, t)$ с~однородными граничными условиями по  координате $z$. 
При этом дифференциальные уравнения (\ref{sl:eq8}), граничные условия (\ref{sl:eq9}) становятся неоднородными с правыми частями $B_1^*$, $B_2^*$, $-H_1$, $-H_2$, $H_3$, $H_4 $, а~в~начальных условиях (\ref{sl:eq11}) вместо $U_0$, $\dot {U}_0$, $W_0$, $\dot {W}_0 $ следует брать $u_0$, $\dot {u}_0 $, $w_0$, $\dot {w}_0 $:
\begin{multline*}
B_1^*  =B_1 -\nabla ^2H_1 +\frac{C_{11} }{C_{33} }\frac{H_1 }{r^2}-
\frac{C_{55} }{C_{33} }\frac{\partial ^2H_1 }{\partial z^2}- \\-
\frac{C_{13} +C_{55}}{C_{33} }\frac{\partial ^2H_2 }{\partial r\partial z}-
\frac{C_{13} -C_{12}}{C_{33} }\frac{1}{r}\frac{\partial H_2 }{\partial z}+
\frac{\partial ^2H_1 }{\partial t^2}, 
\end{multline*}
\begin{multline*}
B_2^*  =B_2 -\frac{C_{13} +C_{55}}{C_{33} }\frac{\partial ^2H_1 }{\partial r\partial z}-
\frac{C_{12} +C_{55} }{C_{33} }\frac{1}{r}\frac{\partial H_1 }{\partial z}- \\-
\frac{C_{55} }{C_{33} }\nabla ^2H_2 -\frac{C_{11} }{C_{33} }\frac{\partial ^2H_2 }{\partial z^2}+\frac{\partial ^2H_2 }{\partial t^2};
\end{multline*}
\begin{gather*}
H_3 =-\frac{\partial H_1 }{\partial r}-\frac{C_{13} }{C_{33} }\Bigl(H_1 +\frac{\partial H_2 }{\partial z} \Bigr) , 
\quad
H_4 =-\Bigl( \frac{\partial H_2 }{\partial r}+\frac{\partial H_1 }{\partial z} \Bigr); 
\\
u_0 (r, z)=U_0 (r, z)-H_1 (r, z, 0), \quad \dot {u}_0 (r, z)=\dot {U}_0 (r, z)-\dot {H}_1  (r, z, 0);
\\
w_0 (r, z)=W_0 (r, z)-H_2 (r, z, 0), \quad \dot {w}_0 (r, z)=\dot {W}_0 (r, z)-\dot {H}_2  (r, z, 0).
\end{gather*}
Функции $f_1 (z)$, $f_2(z)$, $\ldots$, $f_6 (z)$ определяются из следующих дифференциальных уравнений:
\begin{equation}
\label{sl:eq24}
f''_1 (z)=f''_2 (z)=0, \quad f'''_3 (z)=f'''_4(z)=f'''_5(z)=f'''_6(z)=0.
\end{equation}

Применяя к краевой задаче относительно $u(r, z, t)$, $w(r, z, t)$ синус- и косинус"=преобразования Фурье с~конечными 
пределами по переменной $z$,  получаем следующую задачу в~пространстве изображений: 
\begin{equation}
\label{sl:eq25-1}
\begin{array}{l}
\displaystyle
\nabla ^2u_s -
\frac{C_{11} }{C_{33} }\frac{u_s }{r^2}-
\frac{C_{55} }{C_{33} }j_n^2 u_s -
\frac{C_{13} +C_{55} }{C_{33} }j_n \frac{\partial w_c }{\partial r}- \\ \displaystyle\hspace{6cm} -
\frac{C_{13} -C_{12}}{C_{33} }j_n \frac{w_c }{r}-
\frac{\partial ^2u_s }{\partial t^2}=B_{1s}, \\
\displaystyle
\frac{C_{55} }{C_{33} }\nabla ^2w_c -
\frac{C_{11} }{C_{33} }j_n^2 w_c +
\frac{C_{13} +C_{55} }{C_{33} }j_n \frac{\partial u_s }{\partial r}+ \\\displaystyle \hspace{6cm}+
\frac{C_{12} +C_{55} }{C_{33} }j_n \frac{u_s }{r}-
\frac{\partial ^2w_c }{\partial t^2}=B_{2c} ,
\end{array}
\end{equation}
где 
$u_s (r, n, t)= {\cal F}^n_s \bigl(  u(r, z, t) \bigr)$,
$u_{0s} (r, n)= {\cal F}^n_s \bigl(  u_0(r, z) \bigr)$,
$\dot u_{0s} (r, n)= {\cal F}^n_s \bigl(  \dot u_0(r, z) )$,
$B_{1s} =  {\cal F}^n_s \bigl(  B^*_1 \bigr)$,
$H_{1s} =  {\cal F}^n_s \bigl(  H_1 \bigr)$,
$H_{3s} =  {\cal F}^n_s \bigl(  H_3 \bigr)$;
$$  
{\cal F}^n_s\bigl( {}\cdot {} \bigr) =
\int_0^L \bigl( {}\cdot {} \bigr)\sin j_n z\,  dz ,\quad j_n ={n\pi }/{L}; \quad n\in \{0\}\cup \mathbb N;
$$
$w_c (r, n, t)= {\cal F}^n_c \bigl(  w(r, z, t) \bigr)$,
$w_{0c} (r, n)= {\cal F}^n_c \bigl(  w_0(r, z) \bigr)$,
$\dot w_{0c} (r, n)= {\cal F}^n_c \bigl(  \dot w_0(r, z) )$,
$B_{2c} =  {\cal F}^n_c \bigl(  B^*_2 \bigr)$,
$H_{2c} =  {\cal F}^n_c \bigl(  H_2 \bigr)$,
$H_{4c} =  {\cal F}^n_c \bigl(  H_4 \bigr)$;
при $r=k$ и $r= 1$ имеем
\begin{equation}
\label{sl:eq25}
\begin{array}{c}
u_s (k, n, t)=-H_{1s} (k, n, t), \quad w_c (k, n, t)=-H_{2c} (k, n, t); \vspace{1mm} \\
\displaystyle
\Bigl[\frac{\partial u_s }{\partial r}+\frac{C_{13} }{C_{33} }(u_s -j_n w_c ) \Bigr]_{r=1} =H_{3s} (1, n, t), 
 \vspace{1mm} \\
\displaystyle
\Bigr[ \frac{\partial w_c }{\partial r}+j_n u_s  \Bigr]_{r=1} =H_{4c} (1, n, t); 
\end{array}
\end{equation}
при $t=0$ "---
\begin{equation}
\label{sl:eq26}
\begin{array}{c}
u_s (r, n, 0)=u_{0s}, \quad \dot {u}_s (r, n, 0)=\dot {u}_{0s} ; \\ w_c (r, n, 0)=w_{0c} , \quad \dot {w}_c (r, n, 0)=\dot {w}_{0c}.
\end{array}
\end{equation}

На втором этапе решения описанную выше процедуру стандартизации проведем для краевой задачи \eqref{sl:eq25-1}--(\ref{sl:eq26}) по координате $r$, используя 
представление
\begin{equation}
\label{sl:eq27}
u_s (r, n, t)=H_{5s} (r, n, t)+U_s(r, n, t), \; \; w_c (r, n, t)=H_{6c} (r, n, t)+W_c (r, n, t), 
\end{equation}
где  
$$H_{5s} (r, n, t)=f_7 (r)H_{1s} (1, n, t) + f_8 (r)H_{3s}(k, n, t),$$
$$H_{6c} (r, n, t)=f_7 (r)H_{2c} (1, n, t) + f_8 (r)H_{4c}(k, n, t).$$

Если теперь подставить (\ref{sl:eq27}) в (\ref{sl:eq25-1})--(\ref{sl:eq26}) и учесть условия 
\begin{equation}
\label{sl:eq28}
f_7 (1)=f_7 ' (1)=f_8 (k)=f_8 (1)=0, \quad -f_7 (k)=f_8 ' (1)=1, 
\end{equation}
то получим новую начально"=краевую задачу относительно $U_s(r, n, t)$, $W_c (r, n, t)$ с~однородными граничными условиями по координате $r$. 
При этом функции $B_{1s}$, $B_{2c}$, $u_{0s}$, $\dot {u}_{0s}$, $w_{0s}$, $\dot {w}_{0s}$ следует заменить на $B_{1s}^*$, $B_{2c}^*$, $U_{0s}$, $\dot {U}_{0s}$, $W_{0s}$, $\dot {W}_{0s}$ по таким формулам:
\begin{multline*}
B_{1s}^*  = B_{1s} -\nabla ^2H_{5s} +
\frac{C_{11} }{C_{33} }\frac{H_{5s}}{r^2}+
\frac{C_{55} }{C_{33} }j_n^2 H_{5s} + \\+
\frac{C_{13} + C_{55}}{C_{33} }j_n H'_{6c}  +
\frac{C_{13} - C_{12}}{C_{33} }j_n \frac{H_{6c} }{r}+\ddot {H}_{5s} , 
\end{multline*}
\begin{multline*}
B_{2c}^*  =B_{2c} -
\frac{C_{55} }{C_{33} }\nabla ^2H_{6c} +
\frac{C_{11} }{C_{33} }j_n^2 H_{6c} - \\-
\frac{C_{13} + C_{55} }{C_{33} }j_n H'_{5s}  -
\frac{C_{12} + C_{55} }{C_{33} }j_n \frac{H_{5s} }{r}+
\ddot {H}_{6c}; 
\end{multline*}
\begin{gather*}
U_{0s} =u_{0c} -H_{5s} (r, n, 0), \quad  \dot {U}_{0s} =\dot  {u}_{0c} -\dot {H}_{5s} (r, n, 0); \\ 
W_{0c} =w_{0c} -H_{6c} (r, n, 0), \quad  \dot {W}_{0c} =\dot  {w}_{0c} -\dot {H}_{6c} (r, n, 0).
\end{gather*}
Функции $f_7 (r)$, $f_8 (r)$ определяются из следующих уравнений:
\begin{equation}
\label{sl:eq29}
f'''_7 (r)=f'''_8 (r)=0.
\end{equation}

Преобразованная начально"=краевая задача (\ref{sl:eq25})--(\ref{sl:eq26}) относительно \linebreak  $U_s (r, n, t)$, $W_c (r, n, t)$ решается c помощью
структурного алгоритма метода КИП~[\citen{sl:c:6}]. 
Введём на сегменте $[k, 1]$ вырожденное КИП с~неизвестными компонентами вектор"=функции ядра  преобразования 
$K_1 (\lambda_{in} , r)$, $ K_2(\lambda_{in} , r)$: 
\begin{gather}
\label{sl:eq30}
G(\lambda_{in}, n, t)=\int_k^1 \bigl[ U_s (r, n, t) K_1 (\lambda_{in}, r)+W_c (r, n, t)K_2(\lambda_{in}, r)\bigr]r\,dr; \\
\notag
U_s (r, n, t)=\sum\limits_{i=1}^\infty G(\lambda _{in} , n, t) K_1 (\lambda_{in}, r) \| K_{in} \|^{-2} , 
\\
\label{sl:eq31}
W_c (r, n, t) =\sum\limits_{i=1}^\infty 
G(\lambda _{in} , n, t) K_2 (\lambda _{in} , r)\|K_{in}\|^{-2} , \\
\notag
\|K_{in}\|^2=\int_k^1 \bigl[K_1^2 (\lambda _{in}, r)+K_2^2 (\lambda_{in}, r)\bigr]r\,dr , 
\end{gather}
где $\lambda _{in}$, $i\in \mathbb N$ "--- положительные параметры.



После использования структурного алгоритма метода КИП, подробно  изложенного в~работах [\citen{sl:c:5, sl:c:6,sl:c:7}], получаем следующее выражение для трансформанты:
\begin{multline}
\label{sl:eq32}
G(\lambda_{in}, n, t)=G_0 (\lambda _{in} , n)\cos \lambda _{in} t+\dot {G}_0 (\lambda _{in} , t)\lambda _{in}^{-1} \sin \lambda _{in} t-
\\
-\lambda _{in}^{-1} \int_0^t F( \lambda _{in} , n, \tau ) \sin \lambda _{in} (t-\tau )d\tau 
\end{multline}
и с~учётом однородных граничных условий вида (\ref{sl:eq25}) "---  однородную краевую задачу для компонент ядра КИП $K_1 (\lambda_{in}, r)$, $K_2(\lambda_{in}, r)$:
\begin{equation}
\label{sl:eq33}
\begin{array}{c}
\displaystyle
\nabla ^2K_1 -
\frac{C_{11} }{C_{33} }\frac{K_1 }{r^2}-
\frac{C_{55} }{C_{33} }j_n^2 K_1 -
\frac{C_{13} +C_{55}}{C_{33} }j_n {K}'_2 +\lambda _{in}^2 K_1 =0, \vspace{1mm} 
\\
\displaystyle
\frac{C_{55} }{C_{33} }\nabla ^2K_2 -
\frac{C_{11} }{C_{33} }j_n^2 K_2 +
\frac{C_{13} +C_{55}}{C_{33} }j_n \Bigl(K'_1 +\frac{K_1 }{r} \Bigr)+\lambda _{in}^2 K_2 =0;
\end{array}
\end{equation}
при $r=k$ и $r= 1$ "---
\begin{equation}
\label{sl:eq34}
\begin{array}{c}
K_1 ( \lambda _{in}, k )=K_2 (\lambda _{in} , k)=0; \vspace{1mm} \\ 
\displaystyle
\Bigl[ K'_1 +\frac{C_{13} }{C_{33} } (K_1 -j_n K_2 ) \Bigr]_{r=1} =0, 
\quad
\Bigl[K'_2 +j_n K_1 \Bigr]_{r=1} =0.
\end{array}
\end{equation}
Здесь 
\begin{gather*}
F(\lambda_{in}, n, t)=\int_k^1 ( B_{1s}^*  K_1 +B_{2c}^*  K_2 )r\,dr ;
\\
G_0 (\lambda _{in} , n)=\int_k^1[U_{0s} K_1 +W_{0c} K_2 ]r\,dr,\quad
\dot {G}_0 (\lambda _{in} , n)=
\int_k^1 [\dot {U}_{0s} K_1 +\dot {W}_{0c} K_2 ]r\,dr .
\end{gather*}


Система (\ref{sl:eq33}) сводится к следующему дифференциальному уравнению \mbox{4-го} порядка относительно функции $K_1 ( \lambda _{in} , r)$: 
\begin{multline}
\label{sl:eq35}
K_1^{IV} +\frac{2}{r}K'''_1 +\Bigl(b_{1in} -\frac{3}{r}\Bigr)K''_1 +
\Bigl( \frac{b_{1in} }{r}+\frac{3}{r^3} \Bigr)K'_1 + \\+
\Bigl(b_{2in} -\frac{b_{1in} }{r^2}-\frac{3}{r^4} \Bigr)K_1 =0, 
\end{multline}
где 
$$
b_{1in} =\lambda _{in}^2 \Bigl(1+\frac{C_{33} }{C_{55} } \Bigr)+
j_n^2 \frac{C_{13} (C_{13} +2C_{55} )-C_{33}^2 }{C_{33} C_{55} },$$ 
$$b_{2in} =\left( \lambda _{in}^2 -j_n^2  \right)
\Bigl( 
\frac{C_{33} }{C_{55} }\lambda _{in}^2 -j_n^2  \Bigr).$$

Для приведения краевой задачи относительно $U_s (r, n, t)$,  $W_c (r, n, t)$ к (\ref{sl:eq35}) используются условия 
\begin{equation}
\label{sl:eq36}
C_{12} =C_{13},  \quad
C_{11} =C_{33},
\end{equation}
первое из которых позволяет получить самосопряженные исходную относительно $U_s (r, n, t)$, $W_c (r, n, t)$ и~преобразованную 
(\ref{sl:eq33}) системы уравнений, а~второе "--- привести (\ref{sl:eq33}) к~(\ref{sl:eq35}).



Дифференциальное уравнение (\ref{sl:eq35}) допускает факторизацию на коммутативные сомножители
\begin{equation}
\label{sl:eq37}
\left[ \frac{d^2}{dr^2}+\frac{1}{r}\frac{d}{dr}+\Bigl( A_{in}^2 -   \frac{1}{r^2} \Bigr) \right]
\left[ \frac{d^2}{dr^2}+\frac{1}{r}\frac{d}{dr}+\Bigl(\mp D_{in}^2 -\frac{1}{r^2} \Bigr) \right]K_1 (\lambda_{in}, r)=0
\end{equation}
при удовлетворении двух основных случаев отношения коэффициентов для расчёта пьезокерамического цилиндра:
\begin{equation}
\label{sl:eq39-1}
1)\; S  j_n <\lambda _{in} <j_n; \qquad  2)\; j_n <\lambda _{in} ,
\end{equation}
где 
$$
A_{in} =\left[ \frac{b_{1in} }{2}+\Bigl(\frac{b_{1in}^2 }{4}-b_{2in}  \Bigr)^{1/2} \right]^{1/2},\quad 
D_{in} =\bigl(A_{in}^2 \mp b_{1in}  \bigr)^{1/2};$$ 
$$
S=\frac{C_{33}^2 -C_{13} (C_{13} +2C_{55} )}{C_{33} (C_{33} +C_{55} )}. $$

Верхний знак в равенстве (\ref{sl:eq37}) соответствует первому варианту отношения коэффициентов (\ref{sl:eq39-1}), нижний "--- второму случаю. 
Кроме того, следует отметить,  что для пьезокерамического материала [\citen{sl:c:3}] постоянная $S$ изменяется в~пределах от 0,6 до 0,8.


Общий интеграл равенства (\ref{sl:eq37}) определяется одним из следующих выражений:
\begin{equation}
\label{sl:eq40-1}
\begin{array}{l}
1)\; K_1 (\lambda_{in}, r)=N_{1in} J_1 \left( {A_{in} r}  \right)+N_{2in} Y_1 \left( {A_{in} r} \right)+ \\
\hspace{6cm} +N_{3in} I_1 \left( {D_{in} r} 
\right)+N_{4in} \tilde {K}_1 \left( {D_{in} r} \right), \vspace{1mm}
\\
2) \; K_1 (\lambda_{in}, r)=N_{1in} J_1 \left( {A_{in} r} \right)+ N_{2in} Y_1 \left( {A_{in} r} \right)+\\
\hspace{6cm} + N_{3in} J_1 \left( {D_{in} r} 
\right)+N_{4in} Y_1 \left( {D_{in} r} \right)
\end{array}
\end{equation}
в зависимости от случая в \eqref{sl:eq39-1}.
Здесь $J_v ( {} \cdot {} )$, $Y_v ( {} \cdot {} )$, $I_v ( {} \cdot {} )$, $\tilde {K}_v ( {} \cdot {} )$ "--- обыкновенные и~модифицированные 
функции Бесселя 1-го и 2-го родов порядка $v$ [\citen{sl:c:8}].

Используя зависимости между $K_1 (\lambda_{in}, r)$ и $K_2 
(\lambda_{in}, r)$, полученные в процессе приведения (\ref{sl:eq33}) 
к (\ref{sl:eq35}), получаем выражения для второй компоненты ядра преобразований:
\begin{equation}
\label{sl:eq41-1}
\begin{array}{l}
1) \; K_2 (\lambda_{in}, r)=b_{3in} \bigl\{ A_{in} (b_{4in} -A_{in}^2 ) [ N_{1in} J_0 (A_{in} r)+
N_{2in} Y_0 (A_{in} r)]+
\\
\hspace{3cm} 
+N_{3in} D_{in} (b_{4in} +D_{in}^2)I_0(D_{in} r)- \\
\hspace{3cm} \hspace{3cm} - N_{4in} D_{in} (b_{6in} +D_{in}^2)\tilde {K}_0 (D_{in} r) \bigr\}, \vspace{1mm}  \\
2)\; K_2 (\lambda_{in}, r)=b_{3in} \bigl\{ A_{in} (b_{4in} -A_{in}^2 ) [N_{1in} J_0 (A_{in} r )+
N_{2in} Y_0 (A_{in} r)]+
\\
\hspace{3cm} 
+D_{in} (b_{4in} -D_{in}^2 )[N_{3in} J_0 ( D_{in} r)+N_{4in} Y_0 (D_{in} r) ]\bigr\}
\end{array}
\end{equation}
в зависимости от случая в \eqref{sl:eq39-1}. Здесь
\[
b_{3in} =-\frac{C_{55} }{j_n (C_{13} +C_{55} )(\lambda _{in}^2 +b_{2in} j_n^2 )}, 
\quad
b_{4in} =\lambda _{in}^2 +j_n^2 \frac{C_{13} (C_{13} +2C_{55} )}{C_{33} C_{55} }.
\]

Подстановка (\ref{sl:eq40-1}), (\ref{sl:eq41-1}) в~граничные условия (\ref{sl:eq34}) формирует однородную систему уравнений относительно постоянных $N_{1in}$, $N_{2in}$, $N_{3in}$, $N_{4in}$. 
Разыскивая её нетривиальное решение, получаем трансцендентное уравнение для вычисления собственных значений $\lambda _{in}$, а~также выражения для этих
постоянных:
\begin{equation}
\label{sl:eq38}
\det [ B_{s\omega }]=0, \quad  s, \omega\in \{1, 2, 3, 4 \},
\end{equation}
где
\[
B_{1\omega }=Z_1 (P_{in} k), \quad B_{2\omega } =P_{in} (b_{6in} -P_{in}^2 )Z_0 (P_{in} k), 
\]
\[
B_{3\omega } =(C_{13} C_{33}^{-1} -1)Z_1 (P_{in} )+ P_{in} \bigl[1-C_{13} C_{33}^{-1} j_n b_{5in} (b_{6in} -P_{in}^2 )\bigr] Z_0 (P_{in} ), 
\]
\[
B_{4\omega } =\left[ {j_n -b_{5in} P_{in}^2 \left( {b_{6in} -P_{in}^2 } 
\right)} \right]Z_1 \left( {P_{in} } \right);
\]
при $\omega =1$ и $\omega= 2$ имеем $P_{in} =A_{in}$, 
при $\omega =3$ и $\omega= 4$ "--- $P_{in} =D_{in}$, 
при $\omega =1$ и $\omega= 3$ "--- $Z_v ( {} \cdot {} )=J_v ( {} \cdot {} )$, 
при $\omega =2$ и $\omega= 4$ "--- $Z_v ( {} \cdot {} )=Y_v ( {} \cdot {} )$.

Без ограничения общности, принимая 
$
N_{1in} =1, 
$
оставшиеся постоянные интегрирования определяются из решения следующей системы неоднородных уравнений:
\begin{equation}
\label{sl:eq39}
\left[ {{\begin{array}{*{20}c}
 {B_{22} }   & {B_{23} }   & {B_{24} }   \\
 {B_{32} }   & {B_{33} }   & {B_{34} }   \\
 {B_{42} }   & {B_{43} }   & {B_{44} }   \\
\end{array} }} \right]\times \left[ {{\begin{array}{*{20}c}
 {N_{2in} }   \\
 {N_{3in} }   \\
 {N_{4in} }   \\
\end{array} }} \right]= -\left[ {{\begin{array}{*{20}c}
 {B_{21} }   \\
 {B_{31} }   \\
 {B_{41} }   \\
\end{array} }} \right].
\end{equation}

\smallskip

\Section{Расчётные соотношения} 
Заключительным этапом исследования является определение функций $f_1 (z)$, $f_2 (z)$, $\ldots$, $f_6(z)$, $f_7 (r), f_8 (r)$, входящих в~представления (\ref{sl:eq22}), (\ref{sl:eq27}). 
Для этой цели воспользуемся дифференциальными уравнениями (\ref{sl:eq24}), (\ref{sl:eq29}) и~соответствующими граничными условиями (\ref{sl:eq23}), 
(29). 
В результате имеем:
\begin{equation}
\label{sl:eq40}
\begin{array}{c}
\displaystyle
f_1 (z)=\Bigl( 1-\frac{z}{L} \Bigr), \quad f_2 (z)=\frac{z}{L}, \quad
f_3 (z)=\frac{C_{13} }{C_{11} }\Bigl( \frac{z^2}{2L}-z+\frac{3L}{8} \Bigr), \vspace{1mm}\\
\displaystyle
f_4 (z)=-\frac{C_{13} }{C_{11} }\Bigl( \frac{z^2}{2L}-\frac{L}{8} \Bigr), \quad
f_5 (z)=\frac{e_{31} }{e_{15} }\Bigl( \frac{z^2}{2L}-z+\frac{L}{4} \Bigr), \vspace{1mm} \\ \displaystyle
f_6 (z)=-\frac{e_{31} }{e_{15} }\Bigl( \frac{z^2}{2L}-\frac{L}{4} \Bigr), 
\quad
f_7 (r)=-( 1-k)^{-2}(r-1)^2, \vspace{1mm} \\ f_8 (r)=(1-k)^{-1}[r^2-(k+1)r+k].
\end{array}
\end{equation}

Последовательно применяя к~трансформанте (\ref{sl:eq32})  формулы обращения КИП (\ref{sl:eq31}), а~затем конечных синус- и косинус"=преобразований Фурье, с~учётом  (\ref{sl:eq22}), (\ref{sl:eq27}) получаем следующие разложения для $U(r, z, t)$, $W(r, z, t)$:
\begin{equation}
\label{sl:eq41}
\begin{array}{l}
\displaystyle
U(r, z, t)=H_1 (r, z, t)+\sum\limits_{n=0}^\infty \Omega_n^{-1} \biggl[ H_{5s} (r, n, t)+ \vspace{1mm}\\
\displaystyle \
\hspace{4cm}
+ \sum\limits_{i=1}^\infty G(\lambda_{in}, n, t) K_1 (\lambda _{in} , r)\|K_{in}\|^{-2}  \biggr] \sin j_n z, \vspace{1mm}\\
W(r, z, t)=H_2 (r, z, t)+\sum\limits_{n=0}^\infty  \Omega_n^{-1}\biggl[ H_{6c} (r, n, t)+ \vspace{1mm}\\
\hspace{4cm}
+ \sum\limits_{i=1}^\infty G(\lambda_{in}, n, t) K_2 (\lambda _{in} , r)\|K_{in}\|^{-2} \biggr] \cos j_n z.
\end{array}
\end{equation}

Равенства (\ref{sl:eq41}) удовлетворяют дифференциальным уравнениям (\ref{sl:eq8}), краевым (\ref{sl:eq9}), (\ref{sl:eq10}) и~начальным (\ref{sl:eq11}) условиям, т.\,е. представляют замкнутое решение  рассматриваемой начально"=краевой задачи теории упругости. 

\smallskip
\Section{Численный анализ результатов}
В~качестве примера рассматривается пьезокерамический цилиндр ($L=1$, $k=0,5$) из состава ЦТС-19 [\citen{sl:c:2}] с~мембранным закреплением его торцов ($U_1=U_2$) при следующем  воздействии на внешней 
радиальной поверхности: $V(z, t)=V(z)\sin \theta t.$

Рассматриваются различные варианты приложения нагрузки (соответствуют номерам кривых на рис.~1--4): 
\begin{itemize}
\item[1)] по всей поверхности $V(z)=V_0 $;
\item[2)] на нижней половине $V(z)=V_0[1-H(z-L/2)]$;
\item[3)] на нижней и~верхней трети $V(z)=V_0 [1+H(z-2L/3)-H(z-L/3)]$ исследуемого элемента. 
\end{itemize}
Здесь $V_0$ "--- амплитудное значение нагрузки, $\theta$ "--- частота вынужденных колебаний ($\theta =0,3\lambda _{11}$, $\lambda _{11} $ "--- собственное 
значение основного тона колебаний цилиндра [\citen{sl:c:6}]), $H( {} \cdot {} )$  "--- функция Хэвисайда соответствующего аргумента.


\input{./00Pic/sl-1.TpX}

На рис. 1, 2 показаны зависимости, характеризующие изменение во времени $t$ радиальной и~аксиальной компонент вектора перемещений $U(1, L /2, t)$, $W(1, L, t)$;  штриховой линией обозначена осциллограмма приложенной нагрузки. 
Следует отметить, что при действии осесимметричной по высоте (кривые 1, 3) гармонической нагрузки возможно использование допущения о~стационарном режиме вынужденных колебаний,  применяемое при исследовании большинства подобных динамических задач  электроупругости. 
Однако во втором случае (кривая 2) рассматриваемые  точки совершают более сложное движение.

\input{./00Pic/sl-2.TpX}


\input{./00Pic/sl-3.TpX}

\input{./00Pic/sl-4.TpX}

На рис. 3, 4 приведены графики, описывающие изменение амплитудных значений $U(1, z, t)$, $W(1, z, t)$ вдоль аксиальной координаты. 
Результаты показывают, что для кривых 1 и~3 наблюдается одинаковая качественная картина, а~численные значения пропорциональны площади, на которую действует электрический потенциал. 
При действии нагрузки на нижней половине цилиндра (кривая 2) наблюдается более сложная зависимость. 
Можно отметить, что эпюра деформаций $U(1, z, t) \sim z$ имеет разные знаки, а~максимальные перемещения $W(1, z, t)$ на верхней части элемента, в~отличие от симметричной по аксиальной координате нагрузки, наблюдаются не в крайних точках, а~при $z=0,6L$.



В заключение следует отметить, что построенный алгоритм позволяет рассчитать коэффициенты электромеханической связи пьезокерамических \linebreak преобразователей. 


\begin{thebibliography}{9}
\RBibitem{sl:c:1}
\by Шульга~Н.\,А., Болкисев~А.\,М. 
\book Колебания пьезоэлектрических тел 
\publaddr Киев 
\publ Наукова думка
\yr 1990
\finalinfo 228~с


\RBibitem{sl:c:2}
\by Партон~В.\,З., Кудрявцев~Б.\,А. 
\book Электромагнитоупругость пьезоэлектрических и~электропроводных тел 
\publaddr М.
\publ Наука
\yr 1988
\finalinfo 470~с

\RBibitem{sl:c:3}
\by Гринченко~В.\,Т., Улитко~А.\,Ф., Шульга~Н.\,А.
\book Механика связанных полей~элементах конструкций 
\publaddr Киев
\publ Наукова думка
\yr 1989
\finalinfo 279~с


\RBibitem{sl:c:4}
\by Сеницкий~Ю.\,Э.
\paper Осесимметричная динамическая задача для короткого толстостенного цилиндра
\inbook Прикладные проблемы прочности и пластичности. Методы решения задач упругости и пластичности
\publaddr Горький
\publ Горьков. гос. ун-т
\yr 1980
\pages 127--135

\RBibitem{sl:c:5}
\by Сеницкий~Ю.\,Э.
\paper К~решению осесимметричной задачи динамики для анизотропного короткого толстостенного цилиндра
\jour Прикладная механика
\yr  1981
\vol 17
\issue 8
\pages 95--100

\RBibitem{sl:c:6}
\by Сеницкий~Ю.\,Э.
\book Исследование упругого деформирования элементов конструкций при динамических воздействиях методом конечных интегральных преобразований 
\publaddr Саратов
\publ Саратов. ун-т
\yr 1985
\finalinfo 173~с

\RBibitem{sl:c:7}
\by Снеддон~И.\,Н. 
\book Преобразования Фурье
\publaddr М.
\publ Иностр. лит.
\yr 1955
\finalinfo 668~с

\RBibitem{sl:c:8}
\by Янке~Е., Эмде~Ф., Леш~Ф. 
\book Специальные функции. Формулы, графики, таблицы 
\publaddr М.
\publ Наука
\yr 1977
\finalinfo 342~с

\end{thebibliography}

\newpage

\chead[\fancyplain{}{\scriptsize \sl S\,h\,l\,j\,a\,k\,h\,i\,n~D.\,A.}]{\fancyplain{}{\scriptsize \sl  Untied Axisymmetris Dynamic Problem of the Return Piezoeffect \dots}}



\makeengtitlem


%\end{document}
